% Generated draft: requires mathematical and source review.
\documentclass[11pt]{article}
\usepackage[margin=25mm]{geometry}
\usepackage{fontspec}
\setmainfont{Latin Modern Roman}
\usepackage{amsmath,amssymb,mathtools,unicode-math}
\setmathfont{Latin Modern Math}
\usepackage{enumitem,needspace,xurl,hyperref}
\setlength{\parindent}{0pt}
\setlength{\parskip}{5pt}
\setlength{\emergencystretch}{3em}
\newcommand{\N}{\mathbb{N}}
\newcommand{\Z}{\mathbb{Z}}
\newcommand{\Q}{\mathbb{Q}}
\newcommand{\R}{\mathbb{R}}
\newcommand{\C}{\mathbb{C}}
\newcommand{\F}{\mathbb{F}}
\newcommand{\A}{\mathbb{A}}
\newcommand{\PP}{\mathbb{P}}
\newcommand{\E}{\mathbb{E}}
\newcommand{\eps}{\varepsilon}
\newcommand{\dd}{\,\mathrm{d}}
\begin{document}
\section*{UnsolvedMath ID: 10400034}
\textbf{Source problem code:} catalogue entry
\subsection{Definitions and mathematical statement}
A rational Vassiliev invariant of order at most \(d\) extends to singular knots by \(v(K_\times)=v(K_+)-v(K_-)\) and vanishes on knots with more than \(d\) transverse double points. The degree-\(d\) space is an associated graded space, not the set of invariants of exactly degree \(d\). A primitive invariant, normalized to vanish on the unknot, is additive under connected sum.

Let \(\mathcal A(S^1;\Q)\) be the graded rational Jacobi-diagram algebra on an oriented circle, with antisymmetry, IHX and STU relations, before imposing framing independence. A diagram has degree half the number of its univalent and trivalent vertices, excluding the circle itself; every graph component is required to meet the circle. Its connected part \(\mathcal A(S^1;\Q)^{(d)}_{\mathrm{conn}}\) is the degree-\(d\) primitive diagram space, represented by diagrams connected off the circle. The source asks to determine
\[
 p_d=\dim_\Q\mathcal A(S^1;\Q)^{(d)}_{\mathrm{conn}}
 \qquad\text{for every degree }d.
\]
It states this equivalently as the dimension problem for primitive Vassiliev invariants. Its table has a nonzero degree-one framing class, so the framed convention must be distinguished from the quotient by the framing-independence relation \(\mathrm{FI}\).

The supplied Table 1 gives the following exact values or lower bounds, reproduced with their qualifications:
\[
\begin{array}{c|rrrrrrrrrrr}
 d&0&1&2&3&4&5&6&7&8&9&10\\\hline
 p_d&0&1&1&1&2&3&5&8&12&18&27\\
 \dim\mathcal A^{(d)}&1&1&2&3&6&10&19&33&60&104&184\\
 \dim(\mathcal A/\mathrm{FI})^{(d)}&1&0&1&1&3&4&9&14&27&44&80
\end{array}
\]
\[
\begin{array}{c|rrrr}
 d&11&12&13&14\\\hline
 p_d&39&55&\geq78&\geq108\\
 \dim\mathcal A^{(d)}&316&548&\geq932&\geq1591\\
 \dim(\mathcal A/\mathrm{FI})^{(d)}&132&232&\geq384&\geq659
\end{array}
\]
These are supplied historical table entries, not newly verified dimensions. In particular, the entries in degrees 13 and 14 are lower bounds and must not be presented as exact computations. The precise knot framing and diagram conventions need confirmation from the surrounding source, but framing independence is not silently imposed on the requested \(p_d\).
\subsection{Short English statement}
Determine every graded dimension of primitive rational finite-type knot invariants, equivalently the connected circle-diagram spaces in the source's framed convention. The historical table includes lower bounds, not exact values, in degrees 13 and 14.
\subsection{Sources}
\begin{itemize}
\item Local UnsolvedMath registry, permanent ID 10400034, source code catalogue entry; no proposer or proposal year supplied; updated 2026-08-06.
\item Ohtsuki, \emph{Problems on invariants of knots and 3-manifolds} (2002), Problem 2.12, PDF page 33, Table 1: \url{https://msp.org/gtm/2002/04/gtm-2002-04-024s.pdf}. Registry access date 2026-07-29; PDF font-structure extraction. Provenance distinguishes 231 numbered bold headings. Current status and redistribution terms require review.
\item Table 1 explicitly credits source references [67, 224], whose bibliographic entries are not supplied. The summary mentions a Grothendieck--Teichm\"uller Lie-algebra connection without an exact further reference and inaccurately suggests exact dimensions through about degree 14.
\end{itemize}
\end{document}
